% page 286 of the facsimile (image p-285 of the scan)
% block 157.5 x 237.7 mm ; left margin 8.0 mm
\pgc{-12.700mm}{0.000mm}{
\l{\cel{5}278}
\l{}
\l{}
\l{\sou{klaso}.\cel{2}I. Kategorio de civitani. - II. Kategorio de adole-}
\l{\cel{3}canti,+konvokata singla-yare por lia instrukteso mili-}
\l{\cel{3}tala. - III. Kategorio de lernanti qui frequentas singla}
\l{\cel{3}grado di lerno-kurso. - IV. Omna kategorio de personi,}
\l{\cel{3}de kozi, distributita segun la difero pri grado o mem}
\l{\cel{3}(nur) pri naturo. - V. Singla ek la granda grupi de ani-}
\l{\cel{3}mali o de vejetanti. - DEFIRS.}
\l{}
\l{\sou{klasifikar}.\cel{2}(\sou{trans.}) I. (\sou{pri la natur-cienci}). Repartisar}
\l{\cel{3}la individui hierarkiale, segun speci, generi, familii,}
\l{\cel{3}serii, klasi, e c. - II. Distributar metodoze. - DEFIS.}
\l{}
\l{\sou{klasika}. Egardata kom konforma a la reguli di la genero, ed}
\l{\cel{3}uzebla kom modelo. - DEFIRS.}
\l{}
\l{\sou{klastika}. (\sou{geol.}) Qua devenas efike da la rodo mekanismala}
\l{\cel{3}de la aquo movanta. (Ex.: la sedimento-roki).}
\l{}
\l{\sou{klaudikar}. (\sou{netrans.}) Marchar, apogante ne-interegale la gam-}
\l{\cel{3}bi sur sulo, sive pro ke onu doloras ye la gambo, sive}
\l{\cel{3}pro ke un gambo esas plu kurta. - F.}
\l{}
\l{\sou{klauno}. I. Protagonisto bufona di la farsi Angla. - II. Ak-}
\l{\cel{3}toro qua, en ula cirki, distraktas, amuzas, la spektanti}
\l{\cel{3}per prodaji di flexebleso, akompanata da joki e moki. -}
\l{\cel{3}DEFIRS.}
\l{}
\l{\sou{klauzo}.\cel{2}Preskriptajo specala di lego, kontrato, testamento.}
\l{\cel{3}- DEFIRS.}
\l{}
\l{\sou{klavo}. Singla ek la tabuleti levera, blanka o nigra, quin la}
\l{\cel{3}fingri abasas por plear orgeno, piano ; (\sou{analoge)}\cel{1}la klavi}
\l{\cel{3}di skribomashino, di komposto-mashino. - DEFRS.}
\l{}
\l{\sou{klavikordo}. (\sou{muziko)}. Instrumento kun sono fixa, klavaro e}
\l{\cel{3}kordi metala quin pinchas plumo-pinti, qua remplasesis da}
\l{\cel{3}la piano. - DEFIS.}
\l{}
\l{\sou{klavikulo}. (anat.)Osto qua funcionas kom apog-arko inter la}
\l{\cel{3}du shultri. - EFIS.}
\l{}
\l{\sou{klefo}. Peco ek fero, quan onu enduktas aden la truo di seruro,}
\l{\cel{3}e per qua onu funcionigas la mekanismo qua apertas o klo-}
\l{\cel{3}zas ta seruro. - F.}
\l{}
\l{\sou{kleistogam}a.\cel{1}\sou{(bot}.) Dicesas pri floro qua nulinstante apertas}
\l{\cel{3}su, e di qua la fekundigo eventas dum olua klozeso.}
\l{}
\l{\sou{klemar}. (\sou{trans}.)Tenar tre strete : (a) per lokizar en spaco}
\l{\cel{3}tre poka (un o plura kozi o personi) ; (b) per interprox-}
\l{\cel{3}imigar per ligilo (plura kozi o parti di kozo) ; (c) per}
\l{\cel{3}lokizar ici tre an iti (personi o kozi). - D.}
}
